\documentclass[11pt]{article}
\usepackage[a4paper,margin=27mm]{geometry}
\usepackage{amsmath,amssymb,amsthm}
\usepackage[T1]{fontenc}
\usepackage{lmodern}
\usepackage[hidelinks]{hyperref}
\newtheorem{theorem}{Theorem}
\newtheorem{lemma}{Lemma}
\setlength{\parindent}{0pt}
\setlength{\parskip}{6pt}
\title{A disproof of conjecture 00000000331\\\large The literal fixed-numerator assertion}
\author{C0ldSmi1e}
\date{October 5, 2026}
\begin{document}
\maketitle

\section{Statement and scope}
The English and Chinese versions of the repository statement assert that
\[
 \sum_{n=1}^{\infty}\frac{\psi(n)}{n}=+\infty
 \quad\Longrightarrow\quad
 \left|\alpha-\frac{a}{n}\right|<\psi(n)
 \text{ for infinitely many }n,
\]
for almost every $\alpha$, with $a$ arbitrary but fixed. Here denominators
are positive integers, $\psi$ is a nonnegative real-valued function, and
``almost every'' refers to Lebesgue measure. Our function is strictly positive,
bounded by $1/2$, and nonincreasing. We prove failure for every fixed real $a$,
so restricting the numerator to rational, integer, or positive integer values
does not repair the assertion. We treat both the real line and a unit interval.

The numerator stays fixed as $n$ varies. The formula contains neither a varying
integer numerator nor an additive inhomogeneous shift. The conclusion below
concerns this displayed formula, not the research problem suggested by its
Duffin--Schaeffer label. No claim about the status of that research problem
is made. The source imposes no condition $\psi(n)\to0$.

\section{Counterexample and complete proof}
Set $\psi(n)=1/4$ for every $n\geq1$. Write
\[
 S_N=\sum_{n=1}^{N}\frac{1}{4n}.
\]
For $j\geq0$, the block $2^j\leq n<2^{j+1}$ contains $2^j$ terms, each at
least $1/(4\cdot2^{j+1})$. Consequently
\[
 S_{2^k-1}\geq \frac{k}{8}\qquad(k\geq0).
\]
Since the partial sums are nondecreasing, this proves $S_N\to+\infty$:
given a real threshold $M$, choose $k$ with $k/8>M$; then every
$N\geq2^k-1$ satisfies $S_N>M$. Thus the required hypothesis is genuine
divergence to positive infinity, not merely failure of summability.

\begin{theorem}
For every fixed $a\in\mathbb R$, each $\alpha\in[1/2,3/4]$ satisfies
$|\alpha-a/n|<1/4$ for only finitely many positive integers $n$.
\end{theorem}
\begin{proof}
Choose a positive integer $N>4|a|$. For every integer $n\geq N$, we have
$|a|/n<1/4$. If $\alpha\geq1/2$, then
\[
 \left|\alpha-\frac an\right|
 \geq \alpha-\frac an
 \geq \alpha-\frac{|a|}{n}
 > \frac12-\frac14=\frac14.
\]
All successful denominators therefore lie in the finite set
$\{1,\ldots,N-1\}$. This proves the assertion simultaneously for all
$\alpha\in[1/2,3/4]$.
\end{proof}

The failure set contains $[1/2,3/4]$, whose Lebesgue measure is $1/4>0$.
Consequently, the infinitely-many conclusion does not hold for almost every
$\alpha$ on $\mathbb R$, or for almost every $\alpha$ in $[0,1]$. The same failure interval lies inside $(0,1)$ and $[0,1)$, so these
endpoint conventions also retain the positive-measure obstruction. This refutes the implication for every fixed numerator $a$,
even with a strictly positive approximation function.

\section{Formal verification}
The accompanying Lean project uses Lean 4.19.0 and Mathlib v4.19.0.
Its mathematical definitions retain the fixed numerator, the radius
$\psi(n)$, positive denominator indices, and Lebesgue measure.
Positive-integer partial sums express the divergence hypothesis as a limit
to positive infinity. The infinitely-many assertion is stated as infinitude
of the set of successful positive denominators.

The formal proof verifies positivity, divergence, and failure on the stated
interval. Two theorems negate the almost-everywhere conclusions:
\texttt{not\_ae\_real} and \texttt{not\_ae\_unit}.
The theorem \texttt{literal\_claim\_false} refutes the implication for
every fixed real numerator. The encoded implication requires a strictly
positive approximation function; the example meets that requirement.
A further theorem covers nonnegative functions. Open and half-open unit
intervals are also formalized.
The elementary harmonic-series argument above and the formal
harmonic-divergence argument establish the same partial-sum limit.

The project has an explicit default proof target. From \texttt{lean/}, run
\begin{verbatim}
lake build
lake env lean -DwarningAsError=true Conjecture331.lean
lake env lean -DwarningAsError=true Check.lean
\end{verbatim}
\texttt{Check.lean} prints the submitted definitions and the types and axiom
dependencies of the authored theorems. The verification records document a
fresh build, strict source replays, and an exhaustive inventory of constants
originating in the submitted proof module, including generated declarations.
No external numerical computation is a premise of the proof. Local formal
verification and independent internal semantic review are distinct from
maintainer acceptance.
\end{document}
